\documentclass{article}
\usepackage{macros}
\title{Notas de aula - MAT01125}
\author{Leonardo Suarez}
\date{\today}

\begin{document}
\maketitle

Notas de aula da disciplina MAT01125 - Cálculo Numérico II, ministrada pela profa. Cibele Aparecida Ladeia no semestre letivo 26/1. Aulas às terças e quintas-feiras, das 15:30 às 17:10, no laboratório A114. Nas semanas 1 a 9 (área 1), será utilizado como material base os capítulos 8 a 11 do livro \emph{Cálculo Numérico - Um Livro Colaborativo (Python)}, elaborado pelo REAMAT/UFRGS, para abordar derivação e integração numérica e equações diferenciais ordinárias. Nas semanas 10 a 18 (área 2), utilizaremos o material elaborado pelo prof. Dagoberto Justo, com foco em equações diferenciais parciais, diferenças finitas e análise de estabilidade.

Como observado anteriormente, o conteúdo programático será dividido em duas áreas. Na primeira, iniciaremos os estudos de derivação numérica, sua forma geral, fórmulas a partir de pontos uniformemente espaçados e erros de arredondamento. Em seguida, veremos integração numérica, incluindo quadratura por interpolação, regras compostas, método de Romberg, métodos adaptativos, quadratura gaussiana e integrais impróprias. Por fim, estudaremos equações diferenciais ordinárias (EDOs), primeiro com foco em problemas de valor inicial, método de Euler, métodos de Runge-Kutta, métodos de múltiplos passos e estabilidade; e depois problemas de valor de contorno, método \emph{shooting}, métodos de diferenças finitas em problemas lineares e não lineares, análise de estabilidade, consistência e convergência. Na segunda área, abordaremos métodos para resolução de equações diferenciais parciais (EDPs) juntamente com análise de estabilidade.

A avaliação será composta por dois trabalhos. A primeira entrega está prevista para o dia 07/05/2026 e a segunda para o dia 02/07/2026. Há ainda uma recuperação no dia 07/07/2026.

\tableofcontents

\section{Derivação numérica}

Antes de pensar em métodos numéricos para resolver problemas com equações diferenciais, é preciso primeiro adaptar os conceitos de derivada e integral para o mundo discreto. Começamos estudando técnicas para aproximar a derivada de funções reais, isto é, dado um conjunto finito de pontos \(\{(x_i,y_i)\}\), onde \(y_i = f(x_i)\) para \(i = 0, 1, \ldots, n\), queremos encontrar uma aproximação para a derivada de \(f\) em um ponto \(x_0\) a partir dos valores de \(f\) nos pontos vizinhos.

\subsection{Diferenças finitas}

Uma estratégia para aproximar a derivada de \(f\) em um ponto \(x_0\) é utilizar \emph{diferenças finitas}, que são expressões da forma \(f(x_0+h)-f(x_0)\), que, quando divididas por \(h\), chamam-se \emph{quociente de diferenças}. A ideia é que, para \(h\) suficientemente pequeno (mas não tão pequeno a ponto de se ter cancelamento catastrófico), o quociente de diferenças pode ser uma boa aproximação para a derivada de \(f\) em \(x_0\). É claro que tomando-se o limite \(h \to 0\), recupera-se o conceito de derivada. Existem diferentes tipos de diferenças finitas, dependendo dos pontos utilizados para calcular a diferença. As mais comuns são as diferenças finitas progressiva, regressiva e central.

\subsubsection{Diferença finita progressiva de ordem 1}

Seja \(f\) uma função diferenciável. A \emph{diferença finita progressiva} (ou \emph{avançada}) de ordem \(1\)

\begin{equation}
    D^+_h f(x_0) \coloneq \frac{f(x_0+h) - f(x_0)}{h} \approx f'(x_0)
\end{equation}
é uma aproximação para \(f'(x_0)\) com erro de ordem \(\mathcal{O}(h)\)\footnote{Se \(f\) é a função a ser estimada, definida sobre um domínio \(X\), e \(g\) a função de comparação, escreve-se \(f(x) = \mathcal{O}(g(x))\) quando existe uma constante \(M>0\) tal que \(\vert f(x) \vert \leq M g(x)\) para todo \(x \in X\).}. De fato, a partir da série de Taylor de \(f\) em torno de \(x_0+h\), temos

\begin{equation}
    f(x_0+h) = f(x_0) + f'(x_0)h + \frac{f''(\xi)}{2}h^2, \quad h>0, \xi \in (x_0, x_0+h).
\end{equation}
Isolando \(f'(x_0)\), obtemos

\begin{equation}
    f'(x_0) = \frac{f(x_0+h) - f(x_0)}{h} - \frac{f''(\xi)}{2}h = D^+_h f(x_0) - \mathcal{O}(h).
\end{equation}

\subsubsection{Diferença finita regressiva de ordem 1}

De forma análoga, podemos realizar a expansão de Taylor de \(f\) em torno de \(x_0-h\)

\begin{equation}
    f'(x_0) = \frac{f(x_0)-f(x_0-h)}{h} + h\frac{f''(\xi)}{2} = D^-_h f(x_0) + \mathcal{O}(h),
\end{equation}
onde \begin{equation}
    D^-_h f(x_0) \coloneq \frac{f(x_0) - f(x_0-h)}{h}
\end{equation}
é a \emph{diferença finita regressiva} (ou \emph{atrasada}) de ordem \(1\) para \(f'(x_0)\), que também possui erro de truncamento de ordem \(h\).

\subsubsection{Diferença finita central de ordem 2}

Fazendo as expansões de Taylor de \(f\) em torno de \(x_0+h\) e \(x_0-h\) até o termo cúbico e somando as expressões, obtemos

\begin{multline}
    f(x_0+h) - f(x_0-h) = 2f'(x_0)h + h^3\left(\frac{f'''(\xi_{+})-f'''(\xi_{-})}{3!}\right), \\
    \xi_{+} \in (x_0, x_0+h), \xi_{-} \in (x_0-h, x_0).
\end{multline}
Isolando \(f'(x_0)\), temos

\begin{equation}
    f'(x_0) = \frac{f(x_0+h) - f(x_0-h)}{2h} - h^2\left(\frac{f'''(\xi_{+})-f'''(\xi_{-})}{3!}\right),
\end{equation}
onde 

\begin{equation}
    D^0_h f(x_0) \coloneq \frac{f(x_0+h) - f(x_0-h)}{2h}
\end{equation}
é a \emph{diferença finita central} de ordem \(2\) para \(f'(x_0)\), que possui erro de truncamento de ordem \(h^2\), ou ordem 2.

\subsection{Diferença finita para derivada segunda}

Se expandirmos \(f\) em torno de \(x_0+h\) e \(x_0-h\) até o termo de quarta ordem, obtemos

\begin{align}
    f(x_0+h) & = f(x_0) + f'(x_0)h + \frac{f''(x_0)}{2}h^2 + \frac{f'''(x_0)}{3!}h^3 + \mathcal{O}(h^4) \\
    f(x_0-h) & = f(x_0) - f'(x_0)h + \frac{f''(x_0)}{2}h^2 - \frac{f'''(x_0)}{3!}h^3 + \mathcal{O}(h^4)
\end{align}
Somando as expressões e rearranjando os termos, fica

\begin{equation}
    f''(x_0) = \frac{f(x_0+h) - 2f(x_0) + f(x_0-h)}{h^2} - \mathcal{O}(h^2),
\end{equation}
o que motiva a definição da \emph{diferença finita central} de ordem \(2\) para a derivada segunda de \(f\) em \(x_0\):

\begin{equation}
    D^2_h f(x_0) \coloneq \frac{f(x_0+h) - 2f(x_0) + f(x_0-h)}{h^2}.
\end{equation}

\subsection{Erros de arredondamento}

\subsection{Obtenção de fórmulas por polinômios interpoladores}

Até agora vimos como obter fórmulas de diferenças finitas a partir da série de Taylor. No entanto, também é possível obter essas fórmulas utilizando polinômios interpoladores. A ideia é construir um polinômio \(p\) que interpola os pontos \((x_k, f(x_k))\), \(k = 0, 1, \ldots, n\) e, em seguida, calcular a derivada de \(p\) em \(x_0\). A vantagem dessa abordagem é que ela pode ser facilmente generalizada para obter fórmulas de diferenças finitas de ordem superior e para pontos não uniformemente espaçados (embora a conta fique mais ``suja'').

Podemos aproximar a derivada de \(f(x)\) em \(x_0\), \(x_1\) ou \(x_2\) usando uma interpolação de Lagrange de grau \(2\) para os pontos \((x_0, f(x_0))\), \((x_1, f(x_1))\) e \((x_2, f(x_2))\):

\begin{multline}
    f(x) = f(x_0) \frac{(x-x_1)(x-x_2)}{(x_0-x_1)(x_0-x_2)} + f(x_1) \frac{(x-x_0)(x-x_2)}{(x_1-x_0)(x_1-x_2)} + f(x_2) \frac{(x-x_0)(x-x_1)}{(x_2-x_0)(x_2-x_1)} + \\
    + \frac{f'''(\xi(x))}{3!}(x-x_0)(x-x_1)(x-x_2).
\end{multline}
onde
\[
    \frac{f'''(\xi(x))}{3!}(x-x_0)(x-x_1)(x-x_2)
\]
é o termo de erro da interpolação de Lagrange, com \(\xi(x) \in [x_0, x_2]\). Tomando a derivada de \(f\), pondo \(x=x_0\) e considerando uma malha uniforme, com \(x_1 = x_0 + h\) e \(x_2 = x_0 + 2h\), obtemos

\begin{equation}
    f'(x_0) = \frac{1}{h} \left[ -\frac{3}{2}f(x_0) + 2f(x_1) - \frac{1}{2}f(x_2) \right] + \mathcal{O}(h^2).
\end{equation}
Analogamente, pondo \(x = x_1\) e \(x = x_2\), obtemos as fórmulas de diferenças finitas para \(f'(x_1)\) e \(f'(x_2)\):

\begin{align}
    f'(x_1) & = \frac{1}{h} \left[ -\frac{1}{2}f(x_0) + 0f(x_1) + \frac{1}{2}f(x_2) \right] + \mathcal{O}(h^2) \\
    f'(x_2) & = \frac{1}{h} \left[ \frac{1}{2}f(x_0) - 2f(x_1) + \frac{3}{2}f(x_2) \right] + \mathcal{O}(h^2)
\end{align}
donde \(f'(x_1)\) é a diferença finita central de ordem \(2\). Para construir fórmulas de diferenças finitas de ordem superior, basta utilizar um polinômio interpolador de grau maior, o que requer mais pontos de avaliação da função \(f\). Ainda, pode-se considerar pontos não uniformemente espaçados, o que torna a conta significativamente mais complexa.

\subsection{Fórmulas de diferenças finitas}

Outra maneira de obter fórmulas de diferenças finitas para derivadas de qualquer ordem e com alta ordem de precisão é a partir de uma combinação linear dos valores de \(f\) nos \(n\) pontos \(x_k\), com o objetivo de escolher os coeficientes da combinação linear de modo que a aproximação de \(f'(x)\) seja exata para polinômios de grau até \(n-1\). Do Teorema de Taylor, sabe-se que uma função \(f\) suficientemente suave pode ser localmente aproximada por um polinômio \(p\) de grau \(n-1\) a menos de um infinitésimo de ordem superior a \(n-1\). Assim, se a regra de derivação for exata para \(p\), então ela será uma boa aproximação para \(f\).

Dado um conjunto de \(n\) pontos \(\{(x_k), f(x_k))\}_{k=1}^{n}\), escrevemos

\begin{equation}
    f'(x) \approx \sum_{k=1}^{n} w_k f(x_k)
\end{equation}
onde \(w_k\) são os coeficientes a serem determinados de forma que a aproximação seja exata para polinômios de grau até \(n-1\). Sendo \(p(x) = \sum_{k=1}^{n} a_k \phi_k(x)\) um polinômio de ordem \(n-1\) que aproxima \(f(x)\), fixamos a base \(\phi_k(x) = x^{k}\) para \(k = 0, 1, 2, \ldots, n-1\). Fazendo a regra de derivação ser exata para cada polinômio da base, temos o sistema de equações

\begin{equation}
    \begin{bmatrix}
        1 & 1 & \cdots & 1 \\
        x_1 & x_2 & \cdots & x_n \\
        x_1^2 & x_2^2 & \cdots & x_n^2 \\
        \vdots & \vdots & \ddots & \vdots \\
        x_1^{n-1} & x_2^{n-1} & \cdots & x_n^{n-1}
    \end{bmatrix}
    \begin{bmatrix}
        w_1 \\ w_2 \\ w_3 \\ \vdots \\ w_n
    \end{bmatrix}
    =
    \begin{bmatrix}
        0 \\ 1 \\ 2x \\ \vdots \\ (n-1)x^{n-2}
    \end{bmatrix}
\end{equation}
A matriz à esquerda é uma \emph{matriz de Vandermonde}, que é notoriamente mal-condicionada, isto é, conforme o número de pontos \(n\) aumenta, pequenas variações ou erros de arredondamento nos dados de entrada podem causar erros significativos na solução dos pesos \(w_i\).

\subsection{Derivada via ajuste ou interpolação}

\section{Integração numérica}

Uma aproximação de uma integral definida

\begin{equation}
    \int_{a}^{b} f(x) dx \approx \sum_{i=0}^{n} w_i f(x_i)
\end{equation}
é chamada de \emph{regra de quadratura}, onde \(x_i\) são os pontos de avaliação e \(w_i\) são os pesos associados a cada ponto. O objetivo é escolher os pontos \(x_i\) e os pesos \(w_i\) de forma a minimizar o erro da aproximação.

\subsection{Somas de Riemann}

A maneira mais simples de aproximar a integral \[I = \int_{a}^{b} f(x) dx\] é utilizando \emph{somas de Riemann}. Para isso, dividimos o intervalo \([a, b]\) em \(n\) subintervalos de largura \(h = (b-a)/n\) e escolhemos um ponto \(x_i\) em cada subintervalo para avaliar a função \(f\). As somas de Riemann podem ser classificadas em três tipos principais, dependendo da escolha dos pontos \(x_i\):

\begin{align}
    S & = \sum_{i=1}^{n-1} f(x_i)h \qquad\text{(soma de Riemann à esquerda)} \\
    S & = \sum_{i=1}^{n-1} f(x_{i+1})h \qquad \text{(soma de Riemann à direita)} \\
    S & = \sum_{i=1}^{n-1} f\left(\frac{x_i + x_{i+1}}{2}\right)h \qquad \text{(soma de Riemann central)}
\end{align}

\subsection{Regras de Newton-Cotes}

Dado um conjunto de \(n+1\) pontos \(a \leq x_0 < x_1 < \cdots < x_n \leq b\) igualmente espaçados em que o valor de \(f: [a,b] \to \R\) é conhecido, definimos \emph{regras de Newton-Cotes} como

\begin{equation}
    \int_{a}^{b} f(x) dx \approx \sum_{i=0}^{n} w_i f(x_i)
\end{equation}
onde para uma fórmula fechada, tem-se \(x_i = a + ih\), com \(h = (b-a)/n\) e, para uma fórmula aberta, tem-se \(x_i = a + (i+1)h\), com \(h = (b-a)/(n+2)\). O número \(h\) é a \emph{largura do passo}, \(w_i\) são os \emph{pesos}, que podem ser computados como a integral dos polinômios da base de Lagrange. Com efeito, seja \(p_{n-1}\) o polinômio interpolador de Lagrange para os pontos \(\{(x_i, f(x_i))\}_{i=0}^{n}\). Então temos

\begin{align}
    f(x) &= p_{n-1}(x) + R_n(x) \\
    & = \sum_{i=0}^{n} f(x_i) l_i(x) + R_n(x)
\end{align}
onde \(R_n(x)\) é o termo de erro da interpolação de Lagrange, dado por
\begin{equation}
    R_n(x) = \frac{f^{(n)}(\xi)}{n!} \prod_{i=0}^{n} (x - x_i), \quad \xi \in [a, b].
\end{equation}

Substituindo a expressão de \(f(x)\) na integral, obtemos a fórmula de quadratura
\begin{equation}
    \int_{a}^{b} f(x) dx = \sum_{i=0}^{n} f(x_i) \int_{a}^{b} l_i(x) dx + \int_{a}^{b} R_n(x) dx
\end{equation}
que pode ser escrita como

\begin{equation}
    \int_{a}^{b} f(x) dx = \sum_{i=0}^{n} w_i f(x_i) + \int_{a}^{b} R_n(x) dx
\end{equation}
com \(w_i = \int_{a}^{b} l_i(x) dx\) para \(i = 0, 1, \ldots, n\).

A \emph{regra do trapézio} é um caso particular de regra de Newton-Cotes, onde \(n = 1\) e os pontos são \(x_0 = a\) e \(x_1 = b\), com \(h = b-a\). O polinômio interpolador, uma reta, é dado por \(p_1(x) = f(a)l_0(x) + f(b)l_1(x)\), onde \(l_0(x) = \frac{(x-b)}{(a-b)}\) e \(l_1(x) = \frac{(x-a)}{(b-a)}\). Assim, os pesos são \(w_0 = \int_{a}^{b} l_0(x) dx = h/2\) e \(w_1 = \int_{a}^{b} l_1(x) dx = h/2\), resultando na fórmula de quadratura
\begin{equation}
    \int_{a}^{b} f(x) dx \approx \frac{h}{2} [f(a) + f(b)].
\end{equation}

Integrando o erro da interpolação, obtemos a estimativa de erro da regra do trapézio
\begin{equation}
    \int_{a}^{b} R_1(x) dx = -\frac{h^3}{12} f''(\xi), \quad \xi \in (a, b),
\end{equation}
ou seja, a regra do trapézio possui erro de ordem \(\mathcal{O}(h^3)\).

Na \emph{regra de Simpson}, aproximamos \(f\) por uma parábola que interpola os pontos \((a, f(a))\), \((m, f(m))\) e \((b, f(b))\), onde \(m = (a+b)/2\). O polinômio interpolador é dado por
\begin{equation}
    p_2(x) = f(a)l_0(x) + f(m)l_1(x) + f(b)l_2(x),
\end{equation}
onde \(l_0(x) = \frac{(x-m)(x-b)}{(a-m)(a-b)}\), \(l_1(x) = \frac{(x-a)(x-b)}{(m-a)(m-b)}\) e \(l_2(x) = \frac{(x-a)(x-m)}{(b-a)(b-m)}\). Os pesos podem ser calculados fazendo uma translação para a origem (subtraindo \(x_1\) de \(x_2\) e \(x_3\)), donde vem \(w_0 = \int_{a}^{b} l_0(x) dx = h/6\), \(w_1 = \int_{a}^{b} l_1(x) dx = 4h/6\) e \(w_2 = \int_{a}^{b} l_2(x) dx = h/6\). A fórmula de quadratura fica

\begin{equation}
    \int_{a}^{b} f(x) dx \approx \frac{h}{6} [f(a) + 4f(m) + f(b)].
\end{equation}

De forma análoga ao caso do trapézio, encontramos que a regra de Simpson possui erro de ordem \(\mathcal{O}(h^5)\).

\subsection{Obtenção das regras de quadratura}

Da mesma forma que encontramos fórmulas de diferenças finitas a partir de uma combinação linear dos valores de \(f\) nos pontos \(x_k\), também podemos encontrar regras de quadratura a partir de uma combinação linear dos valores de \(f\) nos pontos \(t_i\). Fazendo a regra de quadratura ser exata para a base \(\phi_k(t) = t^{k-1}\), \(1 \leq k \leq n\), ela é exata para o espaço de polinômios de grau até \(n-1\), o que garante uma boa aproximação para funções suficientemente suaves.

Dados \(n\) pontos \([t_1, t_2, \ldots, t_n]\), queremos aproximar

\begin{equation}
    \int_{a}^{b} f(x) dx \approx \sum_{k=1}^{n} w_k f(t_k)
\end{equation}
de modo que a regra de quadratura seja exata para polinômios de ordem \(\leq n-1\). Para isso, aproximamos \(f(t)\) pelo polinômio \(p(t) = \sum_{k=1}^{n} a_k \phi_k(t)\) de ordem \(n-1\). Tomando uma base qualquer, como \(\phi_k(t) = t^{k-1}\), com \(k = 1, 2, \ldots, n\), temos

\begin{equation}
    \int_{a}^{b} \phi_1(t) dt = \int_{a}^{b} 1 dt = b - a = \sum_{k=1}^{n} w_k \phi_1(t_k) = \sum_{k=1}^{n} w_k
\end{equation}
Da mesma forma, para \(k = 2, 3, \ldots, n-1\), temos

\begin{align}
    \int_{a}^{b} \phi_k(t) dt & = \int_{a}^{b} t^{k-1} dt = \frac{b^k - a^k}{k} \\
    & = \sum_{k=1}^{n} w_k \phi_k(t_k) = \sum_{k=1}^{n} w_k t_k^{k-1}
\end{align}
que pode ser escrito na forma matricial como

\begin{equation}
    \begin{bmatrix}
        1 & 1 & \cdots & 1 \\
        t_1 & t_2 & \cdots & t_n \\
        t_1^2 & t_2^2 & \cdots & t_n^2 \\
        \vdots & \vdots & \ddots & \vdots \\
        t_1^{n-1} & t_2^{n-1} & \cdots & t_n^{n-1}
    \end{bmatrix}
    \begin{bmatrix}
        w_1 \\ w_2 \\ w_3 \\ \vdots \\ w_n
    \end{bmatrix}
    =
    \begin{bmatrix}
        b - a \\ (b^2 - a^2)/2 \\ (b^3 - a^3)/3 \\ \vdots \\ (b^n - a^n)/n
    \end{bmatrix}
\end{equation}
Novamente, observa que a matriz à esquerda é uma matriz de Vandermonde.

\subsection{Regras compostas}

Em toda estimativa de erro, o erro depende do tamanho do intervalo de integração. Uma maneira de reduzir o erro é então aumentar o número de subintervalo do intervalo de integração de modo que

\begin{equation}
    \int_{a}^{b} f(x) dx = \sum_{i=1}^{n} \int_{x_i}^{x_{i+1}} f(x) dx
\end{equation}
onde \(a = x_1 < x_2 < \cdots < x_{n+1} = b\) é uma partição do intervalo \([a, b]\) em \(n\) subintervalos. No caso uniforme, tem-se \(x_i = a + (i-1)h\) para \(i = 1, 2, \ldots, n+1\), onde \(h = (b-a)/n\) é a largura de cada subintervalo. Assim, podemos aplicar uma regra de quadratura em cada subintervalo e somar os resultados para obter uma aproximação mais precisa da integral.

\begin{equation}
    \int_{x_i}^{x_{i+1}} f(x) dx \approx \Delta S_i
\end{equation}
e a integral será aproximada por

\begin{equation}
    \int_{a}^{b} f(x) dx \approx \sum_{i=1}^{n} \Delta S_i
\end{equation}
Podemos agora aplicar as regras de quadratura já vistas em cada subintervalo. A \emph{regra composta dos trapézios} toma a forma

\begin{align}
    \int_{a}^{b} f(x) dx &= \sum_{i=1}^{n} \int_{x_i}^{x_{i+1}} f(x) dx \\
    &\approx \sum_{i=1}^{n} \frac{h}{2} [f(x_i) + f(x_{i+1})] \\
    &= \frac{h}{2} \left[f(x_1) + 2\sum_{i=2}^{n} f(x_i) + f(x_{n+1})\right] \\
    &= \frac{h}{2} \left[f(x_1) + f(x_{n+1})\right] + h \sum_{i=2}^{n} f(x_i)
\end{align}

Implementamos o método do trapézio composto em \texttt{Python} para calcular

\begin{equation}
    \int_{a}^{b} f(x) dx = \frac{h}{2} \left[f(x_1) + f(x_{n+1})\right] + h \sum_{i=2}^{n} f(x_i) + \mathcal{O}(h^3).
\end{equation}

\begin{lstlisting}[language=Python, caption={Método do trapézio composto}]
    # escrever codigo
\end{lstlisting}

A \emph{regra composta de Simpson} é dada por

\begin{align}
    \int_{a}^{b} f(x) dx &= \sum_{i=1}^{n} \int_{x_i}^{x_{i+1}} f(x) dx \\
    &\approx \sum_{i=1}^{n} \frac{h}{6} [f(x_i) + 4f\left(\frac{x_i + x_{i+1}}{2}\right) + f(x_{i+1})] \\
    &\approx \frac{h}{3} \left[f(x_1) + 2\sum_{i=1}^{n-1} f(x_{2i+1}) + 4\sum_{i=1}^{n} f(x_{2i}) + f(x_{2n+1})\right] + \mathcal{O}(h^5)
\end{align}
onde \(h = (b-a)/(2n)\), \(x_i = a + (i-1)h\) para \(i = 1, 2, \ldots, 2n+1\) e \(n\) é o número de subintervalos pares.

O código \texttt{Python} a seguir calcula a integral de \(f\) utilizando a regra composta de Simpson.

\begin{lstlisting}[language=Python, caption={Método de Simpson composto}]
    # escrever codigo
\end{lstlisting}

\subsection{Método de Romberg}

O método de Romberg é uma técnica de extrapolação que utiliza as aproximações obtidas por regras de quadratura compostas para melhorar a precisão da estimativa da integral. A ideia é calcular as aproximações \(I(h)\) para diferentes valores de \(h\) e, em seguida, usar essas aproximações para extrapolar o valor da integral quando \(h\) tende a zero.

\begin{thm}
    Se \(f\) é analítica\footnote{Uma função analítica é aquela que pode ser representada por uma série de potências convergente em torno de cada ponto de seu domínio.} em \([a, b]\), então a função \(I(h)\) admite uma representação da forma
    \begin{equation}
        I(h) = I_0 + I_2 h^2 + I_4 h^4 + I_6 h^6 + \cdots
    \end{equation}
\end{thm}
Tomando \(h = (b-a)/2^k\) para \(k = 0, 1, 2, \ldots\), obtemos uma sequência de aproximações \(I_k = I(h)\) que pode ser organizada em uma tabela triangular, onde a primeira coluna contém as aproximações \(I_k\) e as colunas subsequentes contêm as extrapolações obtidas por meio da fórmula de Romberg. A fórmula de extrapolação é dada por

\begin{equation}
    R_{k, j} = \frac{4^j R_{k+1, j-1} - R_{k, j-1}}{4^j - 1}
\end{equation}
onde \(R_{k, 0} = I_k\) para \(k = 0, 1, 2, \ldots\) e \(j = 1, 2, \ldots\).

\subsection{Ordem de precisão}
\subsection{Quadratura de Gauss-Legendre}

Dado um conjunto de \(n\) pontos \(t_1, t_2, \ldots, t_n\) e pesos \(w_1, w_2, \ldots, w_n\), a \emph{quadratura de Gauss-Legendre} é uma regra de quadratura gaussiana que aproxima a integral de uma função \(f\) no intervalo \([-1, 1]\) da seguinte forma:

\begin{equation}
    \int_{-1}^{1} f(x) dx \approx \sum_{i=1}^{n} w_i f(t_i)
\end{equation}
onde os pontos \(t_i\) são as raízes do \emph{polinômio de Legendre de grau \(n\)} e os pesos \(w_i\) são calculados a partir dos pontos \(t_i\) usando a fórmula

\begin{equation}
    w_i = \frac{2}{(1 - t_i^2)[P_n'(t_i)]^2}
\end{equation}
onde \(P_n'(t)\) é a derivada do polinômio de Legendre de grau \(n\). Essa escolha de pontos e pesos garante que a quadratura de Gauss-Legendre seja exata para polinômios de grau até \(2n-1\). Com efeito, considere um polinômio \(p(x)\) de grau \(m \leq 2n-1\). Podemos escrever \(p(x)\) como

\begin{equation}
    p(x) = q(x) P_n(x) + r(x)
\end{equation}
onde \(q(x)\) é um polinômio de grau \(m-n\) e \(r(x)\) é um polinômio de grau menor que \(n\). Ao integrar \(p(x)\) de \(-1\) a \(1\), temos

\begin{equation}
    \int_{-1}^{1} p(x) dx = \int_{-1}^{1} q(x) P_n(x) dx + \int_{-1}^{1} r(x) dx
\end{equation}
A primeira integral é zero, pois \(P_n(x)\) é ortogonal\footnote{Os polinômios de Legendre formam uma base ortogonal para o espaço de funções contínuas no intervalo \([-1, 1]\).} a todos os polinômios de grau menor que \(n\). A segunda integral pode ser calculada exatamente usando as raízes de \(P_n(x)\) como pontos de quadratura.

Para aproximar a integral de uma função \(f\) em um intervalo arbitrário \([a, b]\), podemos realizar uma mudança de variável

\begin{equation}
    x = \frac{b-a}{2} t + \frac{a+b}{2}, \quad t \in [-1, 1]
\end{equation}
tal que

\begin{equation}
    \int_{a}^{b} f(x) dx = \frac{b-a}{2} \int_{-1}^{1} f\left(\frac{b-a}{2} t + \frac{a+b}{2}\right) dt \approx \frac{b-a}{2} \sum_{i=1}^{n} w_i f\left(\frac{b-a}{2} t_i + \frac{a+b}{2}\right).
\end{equation}


\subsection{Integrais impróprias}

\section{Problemas de valor inicial}

\subsection{Rudimentos da teoria de problema de valor inicial}

Uma questão fundamental em matemática aplicada é analisar se uma solução para um problema de valor inicial é útil para predizer o comportamento de um sistema físico ao longo do tempo. Isto é, queremos saber se o problema proposto é \emph{bem posto} (\emph{well-posed}):
\begin{itemize}
    \item Existe solução para o PVI?
    \item A solução é única?
    \item Quão sensível é a solução a pequenas variações nas condições iniciais?
\end{itemize}

\begin{defn}
    Uma função \(f(t,u)\) é \emph{Lipschitz contínua} em um intervalo \(I\) se existe uma constante \(L\) tal que para todo \(t \in I\) e para todo \(u, v \in \R\), tem-se
    \begin{equation}
        \vert f(t,u) - f(t,v) \vert \leq L \vert u - v \vert
    \end{equation}
\end{defn}

\begin{thm} [Picard-Lindelöf]
    Seja \(f(t,u)\) contínua em um retângulo \(R = \{(t,u) : t \in [t_0, t_0 + a], u \in [u_0 - b, u_0 + b]\}\) e Lipschitz contínua em \(u\) para cada \(t\). Então, existe um intervalo \([t_0, t_0 + h]\), com \(h \leq a\), tal que o problema de valor inicial
    \begin{equation}
        \begin{cases}
            u' = f(t,u) \\
            u(t_0) = u_0
        \end{cases}
    \end{equation}
    possui uma única solução \(u(t)\) definida em \([t_0, t_0 + h]\).
\end{thm}

\begin{thm}
    Se \(u(t)\) e \(v(t)\) são soluções de \(u' = f(t,u)\) com \(u(t_0) = u_0\) e \(v(t_0) = v_0\), então
    \begin{equation}
        \vert u(t) - v(t) \vert \leq \vert u_0 - v_0 \vert e^{L(t-t_0)}
    \end{equation}
\end{thm}

\subsection{Método de Euler}

Considere o problema de valor inicial
\begin{equation}
    \begin{cases}
        u' = f(t,u(t)) \\
        u(t_0) = u_0
    \end{cases}
\end{equation}

O \emph{método de Euler} é um método numérico para aproximar a solução \(u(t)\) do problema de valor inicial. A ideia é, partindo do valor inicial \(u(t_0) = u_0\), usar a equação diferencial para obter uma aproximação de \(u(t)\) em um ponto próximo a \(t_0\), e repetir esse processo para obter uma sequência de aproximações. Com o refinamento do passo de discretização, a sequência de aproximações converge para a solução exata \(u(t)\) do problema de valor inicial.

Em outras palavras, se integrarmos ambos os lados da equação diferencial no intervalo \([t_0, t_1]\), obtemos

\begin{equation}
    u(t_1) = u(t_0) + \int_{t_0}^{t_1} f(t, u(t)) dt.
\end{equation}
Aproximando a integral usando o valor de \(f\) em \(t = t_0\), temos a seguinte aproximação \(u_1\) para \(u(t_1)\):

\begin{equation}
    u_1 = u_0 + h f(t_0, u_0)
\end{equation}
Continuando esse processo, podemos obter uma sequência de aproximações \(u_n\) para \(u(t_n)\), onde \(t_n = t_0 + nh\) para \(n = 0, 1, 2, \ldots\). A fórmula de iteração do método de Euler é dada por

\begin{equation}
    u_{n+1} = u_n + h f(t_n, u_n)
\end{equation}
que, juntamente com a condição inicial \(u_0 = u(t_0)\), define o método de Euler.

\subsection{Método de Euler melhorado}

Como visto na subseção anterior, o método de Euler é um método de pouca precisão, o que restringe bastante sua aplicação. Podemos melhorar o método se ao invés de aproximarmos \(f(t,u(t))\) como uma constante igual ao seu valor em \(t=t_n\), aplicarmos a regra do trapézio à integral

\begin{equation}
    u(t_{n+1}) = u(t_n) + \int_{t_n}^{t_{n+1}} f(t, u(t)) dt = u(t_n) + \frac{h}{2} [f(t_n, u(t_n)) + f(t_{n+1}, u(t_{n+1}))] + \mathcal{O}(h^3),
\end{equation}
onde \(h = t_{n+1} - t_n\) é o passo de discretização. Como não conhecemos \(u(t_{n+1})\), podemos usar o método de Euler para obter uma aproximação inicial de \(u(t_{n+1})\) e, em seguida, usar essa aproximação para calcular \(f(t_{n+1}, u_{n+1})\). Então fica

\begin{equation}
    u_{n+1} \approx u_n + \frac{h}{2} [f(t_n, u_n) + f(t_{n+1}, u_n + h f(t_n, u_n))]
\end{equation}
Escrevendo \(k_1 = f(t_n, u_n)\) e \(k_2 = f(t_{n+1}, u_n + hk_1)\) como as derivadas de \(u(t)\) nos pontos \(t_n\) e \(t_{n+1}\), o \emph{método de Euler melhorado} (ou \emph{método de Heun}) assume a forma

\begin{equation}
    \begin{cases}
        u_{n+1} = u_n + \frac{h}{2} (k_1 + k_2) \\
        u_0 = u(t_0)
    \end{cases}
\end{equation}

O erro de truncamento local \(ETL^{n+1}\) do método de Euler melhorado é dado por
\begin{equation}
    ETL^{n+1} = \frac{h^3}{12} u'''(\xi), \quad \xi \in (t_n, t_{n+1}),
\end{equation}
o que mostra que o método de Euler melhorado possui erro de ordem \(\mathcal{O}(h^3)\) e, portanto, é um método de ordem \(2\). Pode ser mostrado que o erro de truncamento global é
\begin{equation}
    ETG^{n+1} = \mathcal{O}(h^2)
\end{equation}

\begin{rem}
    Não se deve concluir que \(ETG = ETL/h\).
\end{rem}

\subsection{Solução de sistemas de equações diferenciais}

Podemos estender as técnicas estudadas para resolver numericamente sistemas de equações diferenciais ordinárias. Considere o sistema de primeira ordem

\begin{equation}
    \begin{cases}
        u_1' = f_1(t, u_1, u_2, \ldots, u_m) \\
        u_2' = f_2(t, u_1, u_2, \ldots, u_m) \\
        \vdots \\
        u_m' = f_m(t, u_1, u_2, \ldots, u_m) \\
        u_1(t_0) = u_{10} \\
        u_2(t_0) = u_{20} \\
        \vdots \\
        u_m(t_0) = u_{m0}
    \end{cases}
\end{equation}
que pode ser escrito de forma compacta como
\begin{equation}
    \begin{cases}
        u'(t) = F(t, u(t)) \\
        u(t_0) = u_0
    \end{cases}
\end{equation}
onde \(u(t) = (u_1(t), u_2(t), \ldots, u_m(t))^T\) é o vetor de soluções e \(F(t, u) = (f_1(t, u), f_2(t, u), \ldots, f_m(t, u))^T\) é o vetor de funções que definem o sistema. Então aplicamos as mesmas fórmulas de iteração do método de Euler para cada componente do vetor \(u(t)\). Fazendo

\begin{equation}
    k_1 = F(t_n, u_n), \quad k_2 = F(t_{n+1}, \tilde{u}_{n+1})
\end{equation}
onde \(\tilde{u}_{n+1} = u_n + hk_1\) é a aproximação de \(u(t_{n+1})\) obtida pelo método de Euler. Assim, o método de Euler melhorado para sistemas de equações diferenciais é dado por
\begin{equation}
    \begin{cases}
        u_{n+1} = u_n + \frac{h}{2} (k_1 + k_2) \\
        u_0 = u(t_0)
    \end{cases}
\end{equation}

\subsection{Solução de equações e sistemas de ordem superior}

É possível estender as técnicas estudadas para resolver numericamente equações diferenciais ordinárias de ordem superior. Considere o problema de segunda ordem dado por:

\begin{equation}
    x''(t) + x'(t) + \sin x(t) = 1,
\end{equation}
sujeito às condições iniciais \(x(0) = 2\) e \(x'(0) = 0\). Para resolver esse problema, podemos reescrevê-lo como um sistema de equações diferenciais de primeira ordem. Definindo \(u_1(t) = x(t)\) e \(u_2(t) = x'(t)\), temos

\begin{equation}
    \begin{cases}
        u_1' = u_2 \\
        u_2' = 1 - u_2 - \sin u_1 \\
        u_1(0) = 2 \\
        u_2(0) = 0
    \end{cases}
\end{equation}
que pode ser escrito de forma compacta como

\begin{equation}
    \begin{cases}
        u' = F(t, u) \\
        u(0) = (2, 0)^T
    \end{cases}
\end{equation}
onde \(u(t) = (u_1(t), u_2(t))^T\) e \(F(t, u) = (u_2, 1 - u_2 - \sin u_1)^T\). Agora podemos aplicar o método de Euler melhorado para obter uma sequência de aproximações \(u^{(n)}\) para \(u(t_n)\), onde \(t_n = nh\) para \(n = 0, 1, 2, \ldots\). Computamos

\begin{align*}
    k_1 &= F(t_n, u^{(n)}) = \begin{bmatrix}
        u_{2}^{(n)} \\
        1 - u_{2}^{(n)} - \sin u_{1}^{(n)}
    \end{bmatrix}
    \\
    \tilde{u}^{(n+1)} &= u^{(n)} + hk_1 = \begin{bmatrix}
        u_{1}^{(n)} + h u_{2}^{(n)} \\
        u_{2}^{(n)} + h (1 - u_{2}^{(n)} - \sin u_{1}^{(n)})
    \end{bmatrix}
    \\
    k_2 &= F(t_{n+1}, \tilde{u}^{(n+1)}) = \begin{bmatrix}
        u_{1}^{(n)} + h u_{2}^{(n)} \\
        1 - (u_{2}^{(n)} + h (1 - u_{2}^{(n)} - \sin u_{1}^{(n)})) - \sin(u_{1}^{(n)} + h u_{2}^{(n)})
    \end{bmatrix}
    \\
\end{align*}
E então, a aproximação \(u^{(n+1)}\) para \(u(t_{n+1})\) é dada por

\begin{equation}
    u^{(n+1)} = u^{(n)} + \frac{h}{2} (k_1 + k_2).
\end{equation}

\subsection{Erro de truncamento}

\begin{defn}
    O \emph{erro de truncamento local} (ETL) é definido como o erro introduzido em cada passo de iteração pelo truncamento da ED supondo conhecida a solução exata no início do intervalo. Um método numérico é de ordem \(p\) se o erro de truncamento local é da ordem \(\mathcal{O}(h^{p+1})\).
\end{defn}

Pode ser mostrado que o método de Euler possui erro de truncamento local \(ETL = \mathcal{O}(h^2)\) e, portanto, é um método de ordem \(1\). Já o método de Euler melhorado possui erro de truncamento local \(ETL = \mathcal{O}(h^3)\) e, portanto, é um método de ordem \(2\).

\begin{defn}
    O \emph{erro de truncamento global} (ETG) é definido como o erro acumulado ao longo de toda a sequência de iterações.
\end{defn}

\subsection{Métodos de Runge-Kutta explícitos}

Os métodos de Runge-Kutta são uma família de métodos numéricos que generalizam o método de Heun para obter métodos de ordem superior. Por ora, estudamos os \emph{métodos de Runge-Kutta explícitos}. Considere o PVI

\begin{equation}
    \begin{cases}
        u' = f(t, u) \\
        u(t_0) = u_0
    \end{cases}
\end{equation}
Vimos que o método de Euler aproxima a integral do lado direito da equação diferencial usando o valor de \(f\) em \(t = t_n\). O método de Heun melhora essa aproximação usando a regra do trapézio, que utiliza os valores de \(f\) em \(t = t_n\) e \(t = t_{n+1}\). Os métodos de Runge-Kutta explícitos admitem pontos intermediários \(\tau_j = t_n + c_j h\) para \(j = 1, 2, \ldots, s\), onde \(0 \leq c_1 \leq c_2 \leq \cdots < c_s \leq 1\) são os nós da quadratura. Assim, aproximamos a integral como

\begin{equation}
    \int_{t_n}^{t_{n+1}} f(t, u(t)) dt \approx h \sum_{j=1}^{s} b_j f(\tau_j, u(\tau_j))
\end{equation}
onde \(b_j\) são os pesos da quadratura. Nota que \(u(\tau_j)\) não é conhecido, mas pode ser aproximado usando os valores de \(f\) nos pontos anteriores \(\tau_1, \tau_2, \ldots, \tau_{j-1}\). Especificamente, podemos escrever

\begin{equation}
    u(\tau_j) \approx \tilde{u}_j = u_n + h \sum_{l=1}^{j-1} a_{jl} f(\tau_l, u(\tau_l))
\end{equation}
e a iteração do método de Runge-Kutta explícito é dada por

\begin{equation}
    u_{n+1} = u_n + h \sum_{j=1}^{s} b_j k_j
\end{equation}
onde \(k_j = f(\tau_j, \tilde{u}_j)\). Os coeficientes \(a_{jl}\), \(b_j\) e \(c_j\) são escolhidos de modo que o método seja de ordem \(p\). Podemos organizar os coeficientes em uma tabela de Butcher, onde a primeira coluna contém os nós \(c_j\), a parte superior direita contém os coeficientes \(a_{jl}\) e a última linha contém os coeficientes \(b_j\):

\begin{center}
    \begin{tabular}{c|ccc}
        \(c_1\) & 0 & \(\cdots\) & 0 \\
        \(c_2\) & \(a_{21}\) & 0 & \(\cdots\) \\
        \(\vdots\) & \(\vdots\) & \(\ddots\) & \(\vdots\) \\
        \(c_s\) & \(a_{s1}\) & \(a_{s2}\) & \(\cdots\) \\
        \hline
        & \(b_1\) & \(b_2\) & \(\cdots\)
    \end{tabular}
\end{center}

\subsection{Métodos de Runge-Kutta implícitos}
\subsection{O método de Taylor}
\subsection{Método de Adams-Bashforth}
\subsection{Método de Adams-Moulton}
\subsection{Método de Adams-Moulton para sistemas lineares}
\subsection{Estratéagia preditor-corretor}
\subsection{Problemas rígidos}
\subsection{Validação e benchmarking}
\subsection{Convergência, consistência e estabilidade}

\section{Problemas de valores de contorno}

\end{document}